% Snapshot of the "Mathematical background" section of main.tex (lines 335-805),
% taken on 2026-11-10 before rewriting that section from scratch.
% Standalone: compile from this folder with `latexmk -pdf math-background_2026-11-10.tex`.
% References to other sections of the survey will show as "??".
\documentclass[11pt]{article}
\usepackage[margin=2.5cm]{geometry}
\usepackage{amsmath,amssymb,amsthm}
\usepackage{enumitem}
\usepackage{xcolor}
\usepackage[numbers,sort&compress]{natbib}
\usepackage{hyperref}
\usepackage{cleveref}

% Theorem-like environments provided by the MDPI class in main.tex
\theoremstyle{plain}
\newtheorem{Theorem}{Theorem}[section]
\newtheorem{Proposition}{Proposition}[section]
\newtheorem{Lemma}{Lemma}[section]
\theoremstyle{definition}
\newtheorem{Definition}{Definition}[section]
\newtheorem{Example}{Example}[section]
\crefname{Theorem}{Theorem}{Theorems}
\crefname{Proposition}{Proposition}{Propositions}
\crefname{Lemma}{Lemma}{Lemmas}
\crefname{Definition}{Definition}{Definitions}
\crefname{Example}{Example}{Examples}

\newcommand{\pnote}[1]{\textcolor{red}{PL: #1}}

\begin{document}

\section{Background: Mathematical Foundations of Symmetry (25/09)}\label{sec:background}
% Group Theory, the algebraic backbone: Group (closure, associativity, identity, inverses), group action (orbits, stabilizers, free vs. transitive actions)
\subsection{Group Theory: The Algebraic Backbone}
\label{sec:group-theory-backbone}

The precise language of symmetry is group theory.
We recall here the minimal apparatus needed: the group axioms, and the theory of group actions, with an emphasis on the orbit--stabilizer picture that recurs throughout.

\subsubsection{Groups}\label{subsec:groups}

\begin{Definition}[Group]\label{def:group}
A \emph{group} is a pair $(G, \cdot)$ consisting of a set $G$ and a binary
operation $\cdot : G \times G \to G$ satisfying:
    \begin{enumerate}[label=(G\arabic*), leftmargin=*, labelsep=0.5em]
        \item \textbf{Closure.} For all $g, h \in G$, $g \cdot h \in G$.
        \item \textbf{Associativity.} For all $g, h, k \in G$,
        \[
            (g \cdot h) \cdot k = g \cdot (h \cdot k).
        \]
        \item \textbf{Identity.} There exists $e \in G$ such that
        \[
            e \cdot g = g \cdot e = g \qquad \text{for all } g \in G.
        \]
        \item \textbf{Inverses.} For every $g \in G$ there exists
        $g^{-1} \in G$ such that
        \[
            g \cdot g^{-1} = g^{-1} \cdot g = e.
        \]
    \end{enumerate}
\end{Definition}

Closure is often stated separately, as above, but is already implicit in the
requirement that $\cdot$ be a map into $G$; we retain it for emphasis, since
in practice (e.g.\ when $G$ arises as a subset of a larger ambient set, such
as permutations of a fixed set or invertible matrices) closure is the axiom
that must be checked first. The identity $e$ is unique, and each inverse
$g^{-1}$ is unique, both by a standard cancellation argument from (G2)--(G4).
If in addition $g \cdot h = h \cdot g$ for all $g,h \in G$, the group is
\emph{abelian}; otherwise it is \emph{non-abelian}.

\begin{Example}
\label{ex:groups}
Groups relevant to the constructions in this survey include:
\begin{itemize}
    \item The \emph{symmetric group} $S_n$ of permutations of
    $\{1, \dots, n\}$ under composition---non-abelian for $n \geq 3$, and the
    natural home for neuron-permutation symmetries of a network layer.
    \item The \emph{general linear group} $GL(n, \mathbb{R})$ of invertible
    $n \times n$ real matrices under multiplication, and its subgroup the
    \emph{orthogonal group} $O(n) = \{ A \in GL(n,\mathbb{R}) : A^\top A = I \}$.
    \item The multiplicative group $\mathbb{R}_{>0}$ of positive reals,
    relevant to positive rescaling symmetries in ReLU networks
    (cf.\ \Cref{sec:parameter-space-symmetries} where such rescalings act
    on a layer's incoming/outgoing weights).
    \item Direct and semidirect products (explained later in this work) of the above, e.g.\
    $S_n \ltimes \mathbb{R}_{>0}^n$, capturing the combined permutation and
    positive-scaling symmetries of a hidden layer.
\end{itemize}
\end{Example}

A \emph{subgroup} $H \leq G$ is a subset closed under $\cdot$ and inverses
and containing $e$; equivalently, $H$ is itself a group under the restriction
of $\cdot$. A map $\varphi : G \to G'$ between groups is a
\emph{homomorphism} if $\varphi(g h) = \varphi(g)\varphi(h)$ for all
$g, h \in G$; it is an \emph{isomorphism} if additionally bijective, in which
case $G$ and $G'$ are structurally indistinguishable as groups.

\subsubsection{Group Actions}
\label{subsec:group-actions}

Groups become tools for organizing geometric or combinatorial structure once
they are allowed to \emph{act} on a set.

\begin{Definition}[Group action]\label{def:group-action}
A (left) \emph{action} of a group $G$ on a set $\mathcal{X}$ is a map
$\alpha : G \times \mathcal{X} \to \mathcal{X}$, written $\alpha(g,x) = g \cdot x$, such that
\begin{enumerate}[label=(A\arabic*), leftmargin=*, labelsep=0.5em]
    \item $e \cdot x = x$ for all $x \in \mathcal{X}$;
    \item $g \cdot (h \cdot x) = (gh) \cdot x$ for all $g, h \in G$,
    $x \in \mathcal{X}$.
\end{enumerate}
Equivalently, an action is a homomorphism $G \to \operatorname{Sym}(\mathcal{X})$, the group
of bijections of $\mathcal{X}$, sending $g \mapsto (x \mapsto g \cdot x)$.
\end{Definition}

Two invariants attached to a point $x \in \mathcal{X}$ organize the geometry of an
action.

\begin{Definition}[Orbit and stabilizer]
\label{def:orbit-stabilizer}
For $x \in \mathcal{X}$, the \emph{orbit} of $x$ under $G$ is
\[
    G \cdot x \;=\; \{\, g \cdot x : g \in G \,\} \subseteq \mathcal{X},
\]
and the \emph{stabilizer} of $x$ is the subgroup
\[
    G_x \;=\; \{\, g \in G : g \cdot x = x \,\} \leq G.
\]
\end{Definition}

Orbits partition $\mathcal{X}$ into equivalence classes under the relation
$x \sim y \iff y \in G \cdot x$, and the quotient set $\mathcal{X} / G$ of orbits is
the natural object classifying $\mathcal{X}$ \emph{up to symmetry}---the set-theoretic
precursor of the weight-space quotients used later to remove permutation and
scaling redundancy from parameter spaces. The stabilizer $G_x$ measures the
symmetry \emph{fixing} a given representative, and the two invariants are
linked by the following basic but central fact.

Before stating the theorem, it is worth isolating the coset construction
that appears in its statement, since it arises itself as the orbit set of a
natural action.

\begin{Example}[Cosets as orbits]
\label{ex:cosets-as-orbits}
Let $H \leq G$ and let $H$ act on the set $\mathcal{X} = G$ via
\[
    \beta(h,x) = x h, \qquad h \in H,\ x \in G.
\]
Composing two such maps gives $\beta(h_2, \beta(h_1,x)) = (x h_1) h_2 = x(h_1 h_2) = \beta(h_1 h_2, x)$,
so this is a \emph{right} action of $H$ on $G$---the composition order is
reversed relative to the left-action convention of
\Cref{def:group-action}. It is transcribed into a left action in
the sense of \Cref{def:group-action} by setting
$h \cdot x := x h^{-1}$ instead, at the cost of an inverse; either
convention produces the same orbits. Taking $\beta$ as written, the orbit of
a point $x \in G$ is
\[
    H \cdot x \;=\; \{\, xh : h \in H \,\} \;=:\; xH,
\]
and the stabilizer is trivial: $xh = x$ forces $h = e$ by cancellation, so
$H_x = \{e\}$ for every $x \in G$---the action is free. The orbit $xH$ is
called the \emph{left coset} of $H$ represented by $x$.
\end{Example}

\begin{Definition}[Left coset]
\label{def:left-coset}
For a subgroup $H \leq G$ and $x \in G$, the \emph{left coset} of $H$
represented by $x$ is
\[
    xH \;:=\; \{\, xh : h \in H \,\} \subseteq G.
\]
The left cosets of $H$ are precisely the orbits of the action in
\Cref{ex:cosets-as-orbits}, and hence partition $G$; the set of all
left cosets is denoted $G/H$. Since that action is free, the map
$h \mapsto xh$ is a bijection $H \xrightarrow{\sim} xH$ for every $x$, so all
left cosets share the cardinality $|H|$. Consequently, when $G$ is finite,
\[
    |G| \;=\; [G:H]\,|H|, \qquad [G:H] := |G/H|,
\]
where $[G:H]$ is the \emph{index} of $H$ in $G$---this is Lagrange's
theorem, an immediate corollary of \Cref{thm:orbit-stabilizer} below
applied to $\mathcal{X} = G$ with $H$ in place of $G$.
\end{Definition}

With this in hand, the set $G/G_x$ appearing in the theorem below is simply
the set of left cosets of the stabilizer subgroup $G_x \leq G$, in the sense
of \Cref{def:left-coset}.

\begin{Theorem}[Orbit--Stabilizer Theorem]
\label{thm:orbit-stabilizer}
Let $G$ act on $\mathcal{X}$ and fix $x \in \mathcal{X}$. The map $g G_x \mapsto g \cdot x$ is a
well-defined bijection between the set of left cosets $G / G_x$ and the
orbit $G \cdot x$. Consequently, if $G$ is finite,
\[
    |G \cdot x| \;=\; [G : G_x] \;=\; \frac{|G|}{|G_x|}.
\]
\end{Theorem}

Two extremal regimes of behavior are of particular interest.

\begin{Definition}[Free and transitive actions]
\label{def:free-transitive}
An action of $G$ on $\mathcal{X}$ is:
\begin{itemize}
    \item \emph{free} if $G_x = \{e\}$ for every $x \in \mathcal{X}$, i.e.\ only the
    identity fixes any point ($g \cdot x = x \implies g = e$);
    \item \emph{transitive} if $\mathcal{X}$ consists of a single orbit, i.e.\ for all
    $x, y \in \mathcal{X}$ there exists $g \in G$ with $g \cdot x = y$.
\end{itemize}
An action that is both free and transitive is \emph{simply transitive} (or
\emph{regular}): for every pair $x, y \in \mathcal{X}$ there is a \emph{unique}
$g \in G$ with $g \cdot x = y$. In this case $\mathcal{X}$ is a $G$-\emph{torsor}, and
the choice of any basepoint $x_0 \in \mathcal{X}$ identifies $\mathcal{X} \cong G$ as $G$-sets via
$g \mapsto g \cdot x_0$---though this identification depends on the
basepoint and carries no group structure of its own on $\mathcal{X}$.
\end{Definition}

Freeness and transitivity are logically independent: an action can be free
without being transitive (e.g.\ $\mathbb{Z}$ acting on itself by translation
has trivial stabilizers but, if $\mathcal{X}$ is taken larger than a single orbit,
fails transitivity), transitive without being free (e.g.\ $S_n$ acting on
$\{1,\dots,n\}$ is transitive but $G_x \cong S_{n-1}$ is nontrivial for
$n \geq 2$), neither, or both. When the action is free, the Orbit--Stabilizer
Theorem collapses to $|G \cdot x| = |G|$: every orbit is a faithful copy of
$G$ itself, with no internal symmetry to quotient by---the situation one
seeks when a group action is meant to model \emph{pure redundancy} (e.g.\
positive rescalings acting freely on the nonzero weights of a ReLU unit) with
no additional accidental coincidences.

Finally, an action is \emph{faithful} if the homomorphism
$G \to \operatorname{Sym}(\mathcal{X})$ of \Cref{def:group-action} is injective,
equivalently $\bigcap_{x \in \mathcal{X}} G_x = \{e\}$; free actions are automatically
faithful, but not conversely. Faithfulness is the natural minimality
condition ensuring $G$ is not acting with redundant, indistinguishable
elements---a standing assumption whenever a symmetry group is identified
\emph{as} a subgroup of $\operatorname{Sym}(\mathcal{X})$ rather than merely mapped into it.




% Discrete (permutation, cyclic...) vs continuous/Lie (Euclidean, Lorentz...) groups

\subsubsection{Discrete versus Continuous (Lie) Groups}
\label{subsec:discrete-vs-lie}

The groups of \Cref{ex:groups} split naturally along a second axis,
orthogonal to the abelian/non-abelian distinction: whether $G$ itself carries
a topology (indeed a smooth manifold structure) compatible with the group
operations, or is merely a set with no such structure to speak of. This
distinction determines whether the tools available are combinatorial
(orbit-counting, coset enumeration) or differential (tangent spaces, flows,
one-parameter subgroups), and both regimes recur throughout this survey.

\subsubsection{Discrete groups}
\label{subsubsec:discrete-groups}

\begin{Definition}[Discrete group]
\label{def:discrete-group}
A group $G$ is \emph{discrete} if it is equipped with the discrete topology
(every subset is open), or, informally, if $G$ is finite or countable with
no accumulation structure relating distinct elements. Equivalently, the only
continuous paths $\gamma : [0,1] \to G$ are the constant ones.
\end{Definition}

\begin{Example}[Discrete groups]
\label{ex:discrete-groups}
\begin{itemize}
    \item The \emph{symmetric group} $S_n$ (\Cref{ex:groups}) is
    finite, of order $n!$, and non-abelian for $n \geq 3$; it acts on
    $\{1,\dots,n\}$ by permutation, and this action is transitive but---as
    already noted---not free.
    \item The \emph{cyclic group} $\mathbb{Z}/n\mathbb{Z}$, the integers
    modulo $n$ under addition, is the prototypical finite abelian group; more
    generally $\mathbb{Z}$ itself, under addition, is an infinite discrete
    abelian group. Every element of $\mathbb{Z}/n\mathbb{Z}$ is a power of a
    single generator, the defining property of \emph{cyclicity}.
    \item Finite subgroups of $O(n)$, such as the dihedral groups
    (symmetries of a regular polygon) or the point groups of crystallography,
    sit discretely inside a continuous ambient group
    (\Cref{subsubsec:lie-groups} below) while remaining discrete as groups
    in their own right.
\end{itemize}
\end{Example}

For discrete $G$, orbits under an action are typically studied by direct
enumeration, and \Cref{thm:orbit-stabilizer} is often used
computationally: $|G\cdot x|$ and $|G_x|$ are counted independently and
cross-checked against $|G|$. This is the regime relevant, e.g., to
neuron-permutation symmetries $S_n \curvearrowright \mathbb{R}^{n}$ acting on
a layer's units.

\subsubsection{Continuous (Lie) groups}
\label{subsubsec:lie-groups}

\begin{Definition}[Lie group]
\label{def:lie-group}
A \emph{Lie group} is a group $G$ that is simultaneously a smooth manifold,
such that the multiplication map $G \times G \to G$, $(g,h) \mapsto gh$, and
the inversion map $G \to G$, $g \mapsto g^{-1}$, are both smooth. Unlike a
discrete group, $G$ then admits nontrivial continuous paths through the
identity, and the tangent space $\mathfrak{g} = T_e G$ at the identity---the
\emph{Lie algebra} of $G$---linearizes the group structure infinitesimally.
\end{Definition}

\begin{Example}[Lie groups]
\label{ex:lie-groups}
\begin{itemize}
    \item The \emph{general linear group} $GL(n,\mathbb{R})$
    (\Cref{ex:groups}) is an open subset of $\mathbb{R}^{n^2}$ under
    the determinant-nonvanishing condition, hence an $n^2$-dimensional Lie
    group; its Lie algebra $\mathfrak{gl}(n,\mathbb{R})$ is all of
    $\mathbb{R}^{n^2}$ under the matrix commutator.
    \item The \emph{Euclidean group} $E(n) = O(n) \ltimes \mathbb{R}^n$
    consists of all distance-preserving transformations of
    $\mathbb{R}^n$---rotations, reflections, and translations---and is the
    Symmetry group of ordinary (flat, positive-definite) Euclidean geometry.
    Its subgroup $SE(n) = SO(n) \ltimes \mathbb{R}^n$ of orientation-preserving
    Isometries are the group of rigid motions.
    \item The \emph{Lorentz group} $O(1,n)$ (or $O(n,1)$, depending on sign
    convention) is the analogue of $O(n)$ for Minkowski space
    $\mathbb{R}^{1,n}$, i.e.\ the group of linear maps preserving the
    indefinite quadratic form $-t^2 + x_1^2 + \cdots + x_n^2$ rather than the
    positive-definite Euclidean one; together with translations, it generates
    the \emph{Poincar\'e group}, the symmetry group of Minkowski spacetime.
    Its noncompactness (boosts are unbounded) already distinguishes its
    representation theory sharply from that of the compact group $O(n)$.
    \item The multiplicative group $\mathbb{R}_{>0}$ and the orthogonal group
    $O(n)$ of \Cref{ex:groups} are themselves Lie groups of dimension
    $1$ and $\binom{n}{2}$ respectively, and their (semidirect) product
    actions on network weight layers---mentioned there in the discrete
    permutation context---are properly continuous, Lie-theoretic actions.
\end{itemize}
\end{Example}

For a Lie group $G$ acting smoothly on a manifold $\mathcal{X}$, orbits
$G \cdot x$ need not be discrete sets but are themselves immersed
submanifolds of $\mathcal{X}$, and \Cref{thm:orbit-stabilizer} upgrades from a
counting statement to a statement about diffeomorphism: when the action is
proper, $G \cdot x \cong G / G_x$ as smooth manifolds, with $G_x \leq G$ now
a closed (possibly positive-dimensional) Lie subgroup rather than a finite
index-computing subgroup. This differential refinement of the
orbit--stabilizer picture is the one in play whenever continuous rescaling
or rotation symmetries---rather than discrete permutation
symmetries---are quotiented out of a parameter space, as in
\Cref{sec:parameter-space-symmetries}.


% Subgroups, direct products, semidirect products

\subsection{Subgroups, Direct Products, and Semidirect Products}
\label{subsec:subgroups-products}

Having distinguished discrete from continuous groups, we now record the
standard constructions for building new groups out of old ones---constructions
already used implicitly above, e.g.\ in writing $E(n) = O(n) \ltimes
\mathbb{R}^n$ or $S_n \ltimes \mathbb{R}_{>0}^n$ for the combined permutation
and rescaling symmetries of a layer.

\subsubsection{Subgroups}
\label{subsubsec:subgroups}

\begin{Definition}[Subgroup, revisited]
\label{def:subgroup}
A subset $H \subseteq G$ is a \emph{subgroup}, written $H \leq G$, if
$e \in H$ and $H$ is closed under the group operation and inversion:
$h_1, h_2 \in H \implies h_1 h_2 \in H$ and $h \in H \implies h^{-1} \in H$.
Equivalently, $H$ is nonempty and $h_1 h_2^{-1} \in H$ for all
$h_1, h_2 \in H$ (the one-step subgroup criterion). Under the restriction of
$\cdot$, $H$ then satisfies (G1)--(G4) of \Cref{def:group} in its
own right.
\end{Definition}

A subgroup $N \leq G$ is \emph{normal}, written $N \trianglelefteq G$, if it
is invariant under conjugation: $g N g^{-1} = N$ for every $g \in G$,
equivalently $g n g^{-1} \in N$ for all $g \in G, n \in N$. Normality is
exactly the condition under which the coset space $G/N$ of
\Cref{def:left-coset} inherits a group structure from $G$, via
$(g_1 N)(g_2 N) := (g_1 g_2) N$---this is well defined precisely because
conjugation cannot move a coset representative out of $N$. The resulting
group $G/N$ is the \emph{quotient group}. Normal subgroups will reappear
below as the fixed piece in a semidirect product, and are the group-theoretic
analogue of the weight-space quotients used to remove symmetry redundancy
elsewhere in this survey.

\begin{Example}
\label{ex:subgroups}
$SO(n) \trianglelefteq O(n)$ (rotations inside all orthogonal
transformations, of index $2$, with $O(n)/SO(n) \cong \mathbb{Z}/2\mathbb{Z}$
recording orientation); the alternating group
$A_n \trianglelefteq S_n$ (even permutations, likewise of index $2$); and the
translation subgroup $\mathbb{R}^n \trianglelefteq E(n)$ inside the Euclidean
group. By contrast, a copy of $O(n) \leq E(n)$ fixing the origin is a
subgroup but \emph{not} normal for $n \geq 1$: conjugating a rotation about
the origin by a nontrivial translation produces a rotation about a different
point, which is not in that particular copy of $O(n)$.
\end{Example}

\subsubsection{Direct products}
\label{subsubsec:direct-products}

\begin{Definition}[Direct product]
\label{def:direct-product}
Given groups $G_1, G_2$, the \emph{direct product} $G_1 \times G_2$ is the
group with underlying set $G_1 \times G_2$ and componentwise operation
\[
    (g_1, g_2) \cdot (g_1', g_2') = (g_1 g_1',\, g_2 g_2').
\]
Both factors sit inside the product as normal subgroups,
$G_1 \times \{e_2\} \trianglelefteq G_1 \times G_2$ and
$\{e_1\} \times G_2 \trianglelefteq G_1 \times G_2$, and elements from the
two factors commute with one another: $(g_1,e_2)(e_1,g_2) = (g_1,g_2) =
(e_1,g_2)(g_1,e_2)$.
\end{Definition}

The direct product is the construction that applies when two symmetries act
\emph{independently}, with no interaction between the factors. It is,
however, too rigid for many of the situations of interest: when one
symmetry's action is itself permuted or twisted by the other, a direct
product cannot express the coupling, and one instead needs the semidirect
product below.

\subsubsection{Semidirect products}
\label{subsubsec:semidirect-products}

Suppose $N$ and $H$ are groups and $H$ acts on $N$ \emph{by automorphisms},
i.e.\ there is a homomorphism $\varphi : H \to \operatorname{Aut}(N)$ (so each
$\varphi(h) : N \to N$ is itself a group isomorphism, not merely a bijection
as in \Cref{def:group-action}).

\begin{Definition}[Semidirect product]
\label{def:semidirect-product}
The \emph{semidirect product} $N \rtimes_\varphi H$ is the group with
underlying set $N \times H$ and operation
\[
    (n_1, h_1) \cdot (n_2, h_2) \;=\; \bigl(n_1 \, \varphi(h_1)(n_2),\; h_1 h_2 \bigr).
\]
This satisfies the group axioms (G1)--(G4) with identity $(e_N, e_H)$ and
inverse $(n,h)^{-1} = (\varphi(h^{-1})(n^{-1}), h^{-1})$; the verification of
(G2) is exactly where the requirement that $\varphi(h)$ be an automorphism
(rather than an arbitrary bijection) is used. The subset
$N \times \{e_H\}$ is a normal subgroup isomorphic to $N$, the subset
$\{e_N\} \times H$ is a subgroup isomorphic to $H$ (normal iff $\varphi$ is
trivial, in which case $N \rtimes_\varphi H$ degenerates to the direct
product $N \times H$), and every element factors uniquely as a product of
one element from each.
\end{Definition}

\begin{Example}[Semidirect products already in use]\label{ex:semidirect-products}
\phantom{u}
\begin{itemize}
    \item \textbf{The Euclidean group.} $E(n) = \mathbb{R}^n \rtimes_\varphi
    O(n)$, where $\varphi : O(n) \to \operatorname{Aut}(\mathbb{R}^n)$ sends a
    matrix $A$ to its defining linear action $v \mapsto Av$. Concretely,
    composing two rigid motions $(v_1, A_1)$ then $(v_2,A_2)$ (translate by
    $v_1$, rotate by $A_1$, then translate by $v_2$, rotate by $A_2$) gives
    $(v_2, A_2)(v_1,A_1) = (v_2 + A_2 v_1,\ A_2 A_1)$---the rotation
    $A_2$ twists the first translation $v_1$ before the two compose, which
    is exactly the semidirect (as opposed to direct) product structure. The
    Poincar\'e group $\mathbb{R}^{1,n} \rtimes O(1,n)$ mentioned in
    \Cref{subsubsec:lie-groups} has the identical shape, with $O(1,n)$ in
    place of $O(n)$.
    \item \textbf{Layer symmetries.} The combined permutation and
    positive-rescaling symmetry of a hidden layer's units,
    $S_n \ltimes \mathbb{R}_{>0}^n$, is the semidirect product with
    $\varphi : S_n \to \operatorname{Aut}(\mathbb{R}_{>0}^n)$ the permutation
    action on coordinates: a permutation $\sigma \in S_n$ acts on a scaling
    tuple $(\lambda_1,\dots,\lambda_n) \in \mathbb{R}_{>0}^n$ by
    $\sigma \cdot (\lambda_i) = (\lambda_{\sigma^{-1}(i)})$, reordering which
    unit gets which scale factor before the scalings themselves compose
    coordinatewise. This is the group whose action on weight space underlies
    the quotient constructions of
    \Cref{sec:parameter-space-symmetries}.
    \item \textbf{Dihedral groups revisited.} $D_n \cong \mathbb{Z}/n\mathbb{Z}
    \rtimes \mathbb{Z}/2\mathbb{Z}$, with the order-$2$ factor acting on the
    cyclic factor by inversion, $\varphi(-1)(k) = -k$; this exhibits the
    discrete symmetry group of \Cref{ex:discrete-groups} itself as a
    semidirect product, showing the construction is not confined to the
    continuous setting.
\end{itemize}
\end{Example}

The direct/semidirect distinction is precisely the question of whether the
two symmetry types present in a problem interact: wherever a rescaling
symmetry and a permutation symmetry (or a rotation and a translation) must be
composed in a fixed order with one twisting the other, the semidirect
product---not the direct product---is the group that actually acts.



% Compactness and Peter-Weyl theorem guaranteeing reducibility, non-compact groups are more difficult: no reducibility of representations, Haar measure

\subsection{Compactness, the Peter--Weyl Theorem, and the Non-Compact Case}
\label{subsec:compactness-peter-weyl}

The discrete/continuous dichotomy of \Cref{subsec:discrete-vs-lie} is not
the only one that matters for representation theory: among Lie groups, a
further split---compact versus non-compact---determines whether
representations behave as tamely as in the finite case or exhibit genuinely
new phenomena. This subsection makes that distinction precise.

\subsubsection{Compact groups}
\label{subsubsec:compact-groups}

\begin{Definition}[Compact group]
\label{def:compact-group}
A Lie group $G$ is \emph{compact} if it is compact as a topological space,
i.e.\ every open cover of $G$ admits a finite subcover, equivalently (since
$G$ is a manifold, hence metrizable) every sequence in $G$ has a convergent
subsequence.
\end{Definition}

\begin{Example}[Compact and non-compact groups already introduced]
\label{ex:compactness}
The orthogonal group $O(n)$ (\Cref{ex:groups}) is compact: it is a
closed and bounded subset of $\mathbb{R}^{n^2}$, boundedness following from
$A^\top A = I \implies \sum_{i,j} A_{ij}^2 = n$. Its subgroup $SO(n)$ is
likewise compact, and finite groups such as $S_n$ or $\mathbb{Z}/n\mathbb{Z}$
are trivially compact. By contrast, $GL(n,\mathbb{R})$, $\mathbb{R}_{>0}$,
the translation group $\mathbb{R}^n \leq E(n)$, and the Lorentz group
$O(1,n)$ (\Cref{subsubsec:lie-groups}) are all non-compact: each contains
sequences with no convergent subsequence (e.g.\ dilations $\lambda \to
\infty$ in $\mathbb{R}_{>0}$, or Lorentz boosts approaching the speed of
light in $O(1,n)$). Consequently $E(n)$ and the Poincar\'e group, being
semidirect products with a non-compact factor, are themselves non-compact,
even though their rotational part $O(n)$, resp.\ $O(1,n)$'s maximal compact
subgroup, is not the source of the non-compactness.
\end{Example}

\subsubsection{Haar measure}
\label{subsubsec:haar-measure}

Representation theory on a Lie group requires a way to integrate functions
on $G$ that is compatible with the group structure.

\begin{Theorem}[Existence of Haar measure]
\label{thm:haar-measure}
On every locally compact topological group $G$ (in particular, every Lie
group) there exists a nonzero, regular Borel measure $\mu$, unique up to
positive scalar, that is \emph{left-invariant}:
\[
    \mu(gS) = \mu(S) \qquad \text{for all } g \in G \text{ and Borel } S \subseteq G.
\]
This $\mu$ is the \emph{(left) Haar measure} of $G$.
\end{Theorem}

In general $\mu$ need not also be right-invariant; groups for which it is
(automatically true for compact and for abelian $G$) are called
\emph{unimodular}. The decisive fact for what follows is:

\begin{Proposition}
\label{prop:compact-haar-finite}
If $G$ is compact, its Haar measure $\mu$ has finite total mass,
$\mu(G) < \infty$, and may be normalized to a left- and right-invariant
probability measure, $\mu(G) = 1$.
\end{Proposition}

If $G$ is non-compact, no such normalization is possible: $\mu(G) = \infty$
necessarily (e.g.\ Lebesgue measure on $\mathbb{R}^n \leq E(n)$, or the
non-compact directions of $O(1,n)$), and integration over all of $G$ ceases
to make sense for a constant function. This single fact is the root cause of
the divergence in behavior described next.

\subsubsection{The Peter--Weyl theorem: compact groups are fully reducible}
\label{subsubsec:peter-weyl}

Fix a finite-dimensional (or, more generally, unitary) representation
$\rho : G \to GL(V)$, i.e.\ a homomorphism into the invertible linear maps of
a vector space $V$ (an action of $G$ on $V$ in the sense of
\Cref{def:group-action}, compatible with the linear structure). A
subspace $W \subseteq V$ is \emph{invariant} if $\rho(g)W \subseteq W$ for
all $g \in G$; $\rho$ is \emph{irreducible} if its only invariant subspaces
are $0$ and $V$, and \emph{completely reducible} if $V$ decomposes as a
direct sum of irreducible invariant subspaces.

\begin{Theorem}[Peter--Weyl]
\label{thm:peter-weyl}
Let $G$ be a compact (Lie) group. Then:
\begin{enumerate}
    \item Every continuous unitary representation of $G$ on a Hilbert space
    decomposes as a direct sum (in general, a Hilbert-space direct sum) of
    finite-dimensional irreducible unitary representations.
    \item Every finite-dimensional representation of $G$ is equivalent to a
    unitary one, and hence---by (1)---is completely reducible.
    \item The matrix coefficients of the irreducible unitary representations
    of $G$ span a dense subspace of $L^2(G,\mu)$ (with $\mu$ the normalized
    Haar measure of \Cref{prop:compact-haar-finite}), giving an
    orthogonal (Peter--Weyl) decomposition
    \[
        L^2(G,\mu) \;\cong\; \bigoplus_{\pi \in \widehat{G}} V_\pi \otimes V_\pi^*,
    \]
    the sum ranging over the set $\widehat{G}$ of (isomorphism classes of)
    irreducible unitary representations of $G$.
\end{enumerate}
\end{Theorem}

The mechanism behind part (2) is the \emph{unitarizability trick}: given any
representation $\rho$ of a compact $G$ on an inner product space $(V,
\langle\cdot,\cdot\rangle)$, average the inner product itself over $G$ using
the finite Haar measure,
\[
    \langle v, w \rangle_G \;:=\; \int_G \langle \rho(g)v,\, \rho(g)w \rangle \, d\mu(g),
\]
which converges precisely because $\mu(G) < \infty$, and is by construction
$G$-invariant, making $\rho$ unitary with respect to $\langle\cdot,\cdot\rangle_G$.
Once $\rho$ is unitary, the orthogonal complement of any invariant subspace
is again invariant, and complete reducibility follows by induction on
$\dim V$. Compactness therefore does not merely simplify representation
theory---it is the single hypothesis that makes the averaging argument, and
hence \Cref{thm:peter-weyl} in its entirety, possible.

\subsubsection{Non-compact groups: no such guarantee}
\label{subsubsec:noncompact-representations}

When $G$ is non-compact, \Cref{prop:compact-haar-finite} fails
($\mu(G) = \infty$), the averaging trick above no longer converges, and both
conclusions of \Cref{thm:peter-weyl} can genuinely fail.

\begin{itemize}
    \item \textbf{Complete reducibility can fail.} A finite-dimensional
    representation of a non-compact $G$ need not be unitarizable, and an
    invariant subspace need not admit an invariant complement. The
    prototypical example is the standard action of $\mathbb{R}$ (as the
    subgroup of upper unitriangular matrices) on $\mathbb{R}^2$ via
    $t \cdot (x,y) = (x + ty,\, y)$: the line $y = 0$ is invariant, but no
    invariant complementary line exists, so this representation is
    reducible yet \emph{not} completely reducible---a phenomenon that
    cannot occur for a compact group by \Cref{thm:peter-weyl}.
    \item \textbf{Irreducible unitary representations are typically
    infinite-dimensional.} For non-compact, non-abelian semisimple groups
    such as the Lorentz group $O(1,n)$, essentially all irreducible unitary
    representations (beyond trivial finite-dimensional ones on which the
    noncompact directions act trivially) are infinite-dimensional, organized
    into families---principal series, complementary series, discrete series
    (when they exist)---rather than a single discrete spectrum $\widehat{G}$
    as in part (3) of \Cref{thm:peter-weyl}. The Peter--Weyl
    decomposition of $L^2(G,\mu)$ itself must be replaced by a
    \emph{direct integral} over these families,
    \[
        L^2(G,\mu) \;\cong\; \int_{\widehat{G}}^{\oplus} V_\pi \otimes V_\pi^* \, d\nu(\pi),
    \]
    with $\nu$ a Plancherel measure on the (now typically continuous)
    unitary dual $\widehat{G}$---the substance of Harish-Chandra's
    representation theory of semisimple Lie groups, and considerably harder
    than the compact case.
\end{itemize}

The practical consequence for constructions elsewhere in this survey is
this: whenever a symmetry group acting on a parameter or feature space is
compact (finite permutation groups, $O(n)$, $SO(n)$), invariant
decompositions and averaged inner products are available essentially for
free via \Cref{thm:peter-weyl}; whenever it is non-compact (scaling
groups $\mathbb{R}_{>0}$, translation groups, the full Euclidean or Lorentz
groups), no such decomposition is automatic, and invariance must instead be
built by hand---e.g.\ by restricting to a compact subgroup, passing to a
homogeneous space with a genuinely invariant (rather than merely
quasi-invariant) measure, or working with unbounded, possibly
infinite-dimensional representations directly.

% Give an example of a simple group, reference a textbook

\subsubsection{Simple groups: an example}
\label{subsubsec:simple-groups}

The normal subgroups introduced in \Cref{subsubsec:subgroups} single out a
distinguished class of groups: those admitting none.

\begin{Definition}[Simple group]
\label{def:simple-group}
A group $G$ is \emph{simple} if $G \neq \{e\}$ and the only normal subgroups
of $G$ are $\{e\}$ and $G$ itself.
\end{Definition}

Simple groups are the group-theoretic analogue of prime numbers: every
finite group can be built from simple groups by a finite sequence of
extensions (a \emph{composition series}), so that classifying finite simple
groups amounts to classifying the irreducible building blocks out of which
all finite groups are assembled. Abelian simple groups are easy to
classify---$\mathbb{Z}/p\mathbb{Z}$ for $p$ prime, since any nontrivial
proper subgroup of an abelian group is automatically normal, so simplicity
forces the group to have no nontrivial proper subgroups at all. The first
genuinely non-abelian example is the following.

\begin{Example}[$A_5$ is simple]
\label{ex:A5-simple}
The alternating group $A_5$ (even permutations of $\{1,\dots,5\}$,
$|A_5| = 60$) introduced alongside $A_n \trianglelefteq S_n$ in
\Cref{ex:subgroups} is the smallest non-abelian simple group. A
standard proof examines the conjugacy classes of $A_5$, of sizes
$1, 15, 20, 12, 12$, and shows that no proper nonempty union of these class
sizes summing to a divisor of $60$ (other than $\{1\}$ and the full set)
exists---since a normal subgroup must be a union of conjugacy classes
(including $\{e\}$) whose total size divides $|A_5|$ by
\Cref{thm:orbit-stabilizer}, this rules out every candidate normal
subgroup besides $\{e\}$ and $A_5$ itself. Geometrically, $A_5$ is also the
rotation group of the icosahedron/dodecahedron, giving a second,
more visual route to the same conclusion via its action on the $5$
inscribed cubes of a dodecahedron.

Simplicity of $A_5$ is what obstructs solving the general quintic by
radicals: Galois theory reduces solvability-by-radicals to the
\emph{solvability} of a certain group (existence of a composition series
with abelian factors), and a simple non-abelian composition factor such as
$A_5$ blocks this, in sharp contrast to $S_n$ for $n \leq 4$, whose
composition series bottoms out entirely in abelian pieces. A detailed proof
of simplicity of $A_5$, together with its role in this Galois-theoretic
obstruction, can be found in Dummit and Foote,
\emph{Abstract Algebra}~\citep{dummit2003}, \S 3.5 and \S 6.4 (in various
editions, in the sections on simple groups and Galois theory,
respectively).
\end{Example}

The classification of \emph{all} finite simple groups---completed over
several decades by dozens of authors and running, across all contributing
papers, to an estimated ten thousand journal pages---identifies $A_5$ as the
smallest member of the infinite family of alternating groups $A_n$
($n \geq 5$), one of eighteen infinite families (together with cyclic groups
of prime order, the alternating groups, and sixteen families of groups of
Lie type) plus twenty-six sporadic exceptions. This classification lies far
outside the scope of the present survey; we record it here only to situate
$A_5$ as the first, and most tractable, instance of a phenomenon that is in
general extraordinarily deep.

%
\subsection{Representation Theory}
\label{subsec:representation-theory}

We now assemble the representation-theoretic vocabulary used implicitly in
\Cref{subsubsec:peter-weyl}--\Cref{subsubsec:noncompact-representations}
into a self-contained toolkit, since it is this vocabulary---not the group
axioms themselves---that is invoked directly wherever equivariance and
invariance are built into a model architecture later in this survey.

\subsubsection{Linear representations}
\label{subsubsec:linear-representations}

\begin{Definition}[Linear representation]
\label{def:linear-representation}
A \emph{linear representation} of a group $G$ on a vector space $V$ (over
$\mathbb{R}$ or $\mathbb{C}$) is a group homomorphism
\[
    \rho : G \longrightarrow GL(V),
\]
i.e.\ a map satisfying $\rho(g_1 g_2) = \rho(g_1)\rho(g_2)$ and
$\rho(e) = \operatorname{id}_V$ for all $g_1, g_2 \in G$---equivalently (as already
noted in \Cref{subsubsec:peter-weyl}), a group action of $G$ on $V$
(\Cref{def:group-action}) through \emph{linear} bijections of $V$
rather than arbitrary bijections of a set. We write $\dim \rho := \dim V$,
and call $\rho$ \emph{unitary} if $V$ carries a $G$-invariant inner product,
i.e.\ $\rho(g) \in U(V)$ for all $g$.
\end{Definition}

Two representations $\rho_1 : G \to GL(V_1)$, $\rho_2 : G \to GL(V_2)$ are
\emph{equivalent} (or isomorphic) if there is a linear isomorphism
$T : V_1 \to V_2$ intertwining the actions,
$T \rho_1(g) = \rho_2(g) T$ for all $g \in G$; such a $T$ is called an
\emph{intertwiner}, or a \emph{$G$-equivariant map}, and the representation
theory of $G$ is, at bottom, the classification of representations up to
this notion of equivalence.

\begin{Example}[Representations already encountered]
\label{ex:representations-recap}
The \emph{defining representation} of $O(n)$ on $\mathbb{R}^n$ (matrices
acting on column vectors, \Cref{ex:groups}) is a linear
representation in this sense; the \emph{regular representation} of a finite
group $G$ on $\mathbb{C}[G] = \bigoplus_{g \in G} \mathbb{C}\, g$, defined by
$\rho(g)\bigl(\sum_h c_h\, h\bigr) = \sum_h c_h\,(gh)$, encodes $G$ acting on
itself by left multiplication (the free action of
\Cref{ex:cosets-as-orbits} with $H = G$, linearized); and the
\emph{trivial representation} $\rho(g) = \operatorname{id}_V$ for all $g$ is
present inside every representation as the subspace of $G$-fixed vectors.
\end{Example}

\subsubsection{Irreducibility and reducibility for compact groups}
\label{subsubsec:irreps-recap}

Irreducibility and complete reducibility were already defined in
\Cref{subsubsec:peter-weyl} in order to state the Peter--Weyl theorem; we
restate them here as the central organizing notions of the subject proper.

\begin{Definition}[Irreducible representation]
\label{def:irrep}
A representation $\rho : G \to GL(V)$ is \emph{irreducible} (an
\emph{irrep}) if $V \neq 0$ and $V$ has no $G$-invariant subspace other than
$0$ and $V$. It is \emph{completely reducible} if $V$ is a direct sum of
irreducible invariant subspaces.
\end{Definition}

For $G$ compact, \Cref{thm:peter-weyl} guarantees that every
finite-dimensional representation is completely reducible, and organizes the
whole of $L^2(G)$ into the discrete direct sum indexed by the unitary dual
$\widehat{G}$; this is the structural fact that recent work on
group-equivariant deep learning leans on directly, decomposing feature
spaces into isotypic (irrep-indexed) components in order to constrain
network layers to be exactly $G$-equivariant by design rather than by
regularization---see, e.g., the treatment of spherical CNNs and Clebsch--Gordan
decompositions in \citep{kondor2018p3icml} and the general group-equivariant
convolutional framework of \citep{cohen2014}, both of which exploit
compactness of the relevant symmetry group ($SO(3)$, or finite groups such as
the rotation/reflection group of a square lattice, respectively) precisely
to invoke the Peter--Weyl decomposition of \Cref{subsubsec:peter-weyl}. As
emphasized in \Cref{subsubsec:noncompact-representations}, this guarantee
evaporates once $G$ is non-compact, which is exactly why equivariant
architectures for non-compact symmetry groups (translations, positive
rescalings) require separate, ad hoc constructions rather than a single
unified decomposition.

\subsubsection{Tensor algebra}
\label{subsubsec:tensor-algebra}

Representations combine via the usual multilinear-algebra operations, each
of which is again a representation.

\begin{Definition}[Operations on representations]
\label{def:tensor-operations}
Given representations $\rho_1 : G \to GL(V_1)$ and $\rho_2 : G \to GL(V_2)$:
\begin{itemize}
    \item The \emph{direct sum} $\rho_1 \oplus \rho_2$ acts on
    $V_1 \oplus V_2$ by $(\rho_1 \oplus \rho_2)(g)(v_1,v_2) =
    (\rho_1(g)v_1,\, \rho_2(g)v_2)$.
    \item The \emph{tensor product} $\rho_1 \otimes \rho_2$ acts on
    $V_1 \otimes V_2$ by
    $(\rho_1 \otimes \rho_2)(g) = \rho_1(g) \otimes \rho_2(g)$, i.e.\
    $(\rho_1\otimes\rho_2)(g)(v_1 \otimes v_2) = \rho_1(g)v_1 \otimes \rho_2(g)v_2$,
    extended linearly.
    \item The \emph{dual representation} $\rho_1^*$ acts on
    $V_1^* = \operatorname{Hom}(V_1,\mathbb{C})$ by
    $(\rho_1^*(g)\varphi)(v) = \varphi(\rho_1(g)^{-1} v)$, the inverse
    ensuring $\rho_1^*$ is itself a homomorphism rather than an
    anti-homomorphism.
    \item Iterating the tensor product gives the full \emph{tensor algebra}
    $T(V) = \bigoplus_{k \geq 0} V^{\otimes k}$ of a representation
    $\rho: G \to GL(V)$, on which $G$ acts on each graded piece
    $V^{\otimes k}$ via $\rho^{\otimes k}$; its symmetric and antisymmetric
    quotients $\operatorname{Sym}^k(V)$ and $\Lambda^k(V)$ are likewise
    $G$-representations, obtained by restricting $\rho^{\otimes k}$ to the
    (anti)symmetric subspace of $V^{\otimes k}$.
\end{itemize}
\end{Definition}

Unlike direct sum, the tensor product of two irreducibles is typically
\emph{not} irreducible: for $G$ compact, complete reducibility
(\Cref{thm:peter-weyl}) guarantees $\rho_1 \otimes \rho_2$ decomposes
as a direct sum of irreps,
\[
    \rho_1 \otimes \rho_2 \;\cong\; \bigoplus_{\pi \in \widehat{G}} N_{\pi}^{\rho_1,\rho_2}\; \pi,
\]
with nonnegative integer multiplicities $N_\pi^{\rho_1,\rho_2}$; determining
these multiplicities explicitly (the \emph{Clebsch--Gordan problem}, e.g.\
for $G = SO(3)$ as used in \citep{kondor2018p3icml}) is itself a substantial
part of the representation theory of any given compact group.

\subsubsection{Schur's lemma}
\label{subsubsec:schurs-lemma}

The following elementary but pivotal fact controls exactly which
intertwiners between irreducibles can exist, and underlies both the
uniqueness statement in Peter--Weyl and the practical computation of
Clebsch--Gordan multiplicities above.

\begin{Lemma}[Schur]
\label{lem:schur}
Let $\rho_1 : G \to GL(V_1)$ and $\rho_2 : G \to GL(V_2)$ be irreducible
representations over $\mathbb{C}$, and let $T : V_1 \to V_2$ be an
intertwiner, $T\rho_1(g) = \rho_2(g) T$ for all $g \in G$. Then:
\begin{enumerate}
    \item Either $T = 0$, or $T$ is an isomorphism (so $\rho_1 \cong \rho_2$).
    \item If $\rho_1 = \rho_2 = \rho$ (so $T : V \to V$), then
    $T = \lambda\, \operatorname{id}_V$ for some scalar $\lambda \in \mathbb{C}$.
\end{enumerate}
\end{Lemma}

\begin{proof}[Proof sketch]
For (1): $\ker T \subseteq V_1$ is a $G$-invariant subspace, since
$\rho_1(g)v \in \ker T$ whenever $v \in \ker T$ (as $T\rho_1(g)v =
\rho_2(g)Tv = 0$); irreducibility of $\rho_1$ forces $\ker T \in \{0, V_1\}$.
Symmetrically, $\operatorname{im}\, T \subseteq V_2$ is $G$-invariant, so
irreducibility of $\rho_2$ forces $\operatorname{im}\, T \in \{0, V_2\}$. Excluding
the trivial map leaves $\ker T = 0$ and $\operatorname{im}\, T = V_2$, i.e.\ $T$ is
an isomorphism. For (2): since $\mathbb{C}$ is algebraically closed, $T$ has
an eigenvalue $\lambda$; then $T - \lambda\,\operatorname{id}_V$ is again an
intertwiner of $\rho$ with itself, but now with nontrivial kernel, so by (1)
it must be the zero map, giving $T = \lambda\,\operatorname{id}_V$.
\end{proof}

Part (2) is the source of the diagonal, isotypic block structure exploited
whenever a $G$-equivariant linear layer is constructed for compact $G$: any
linear map commuting with the $G$-action on an isotypic component (a direct
sum of copies of one irrep $\pi$) is forced by Schur's lemma to act as a
single scalar (or, more generally, a small matrix indexed by the
multiplicity space) on that component, which is precisely the parameter
reduction---from a dense weight matrix to a handful of learnable
scalars per irrep---underlying the equivariant convolution constructions of
\citep{cohen2014} and the tensor-field/Clebsch--Gordan layers of
\citep{kondor2018p3icml}.

% Isotypic decomposition and its role in constraining architecture: used in weiler2019anips.md, xu2022p3icml.md, yarotsky2018r.md, gibson2024.md

\subsubsection{Isotypic decomposition and its role in constraining architecture}\label{sec:isotypic-decomposition}

The Peter--Weyl and Clebsch--Gordan decompositions of
\Cref{subsubsec:irreps-recap}--\Cref{subsubsec:tensor-algebra} organize a
representation not merely into individual irreps, but into blocks of
\emph{repeated} irreps, and it is this coarser block structure---rather than
the finer decomposition into individual irreducible summands, which is
unique only up to the choice of an arbitrary basis within each block---that
is the object directly constrained in equivariant network design.

\begin{Definition}[Isotypic decomposition]
\label{def:isotypic-decomposition}
Let $\rho : G \to GL(V)$ be a representation of a compact group $G$, and let
$\widehat{G}$ denote its unitary dual (\Cref{thm:peter-weyl}). For
$\pi \in \widehat{G}$, the \emph{$\pi$-isotypic component} of $V$ is the sum
of all irreducible invariant subspaces of $V$ isomorphic to $\pi$; it is
canonically of the form $V_\pi \cong \pi^{\oplus m_\pi} \cong \pi \otimes
\mathbb{C}^{m_\pi}$, where $m_\pi \in \mathbb{Z}_{\geq 0}$ is the
\emph{multiplicity} of $\pi$ in $\rho$, and $\mathbb{C}^{m_\pi}$ (the
\emph{multiplicity space}) carries no $G$-action at all. The
\emph{isotypic decomposition}
\[
    V \;=\; \bigoplus_{\pi \in \widehat{G}} V_\pi \;=\; \bigoplus_{\pi \in \widehat{G}} \pi \otimes \mathbb{C}^{m_\pi}
\]
is then canonical (unlike the decomposition of each $V_\pi$ into individual
irreducible lines, which depends on a basis choice within
$\mathbb{C}^{m_\pi}$), since distinct isotypic components are automatically
orthogonal and $G$-invariant, and the summands within a single $V_\pi$ are
not individually distinguished by the $G$-action.
\end{Definition}

\begin{Proposition}[Schur blocks intertwiners]
\label{prop:schur-blocks}
Let $V = \bigoplus_\pi \pi \otimes \mathbb{C}^{m_\pi}$ and
$V' = \bigoplus_\pi \pi \otimes \mathbb{C}^{m_\pi'}$ be isotypic
decompositions of two representations of $G$. Then any $G$-equivariant
linear map $T : V \to V'$ (\Cref{def:linear-representation})
vanishes between non-matching isotypic components, and restricts on each
matching component to a map of the form
\[
    T\big|_{V_\pi} \;=\; \operatorname{id}_\pi \otimes A_\pi, \qquad A_\pi \in \operatorname{Hom}(\mathbb{C}^{m_\pi}, \mathbb{C}^{m_\pi'}),
\]
i.e.\ $T$ acts as the identity on the irrep factor and as an
\emph{arbitrary} linear map $A_\pi$ on the multiplicity space alone.
\end{Proposition}

This is precisely Schur's Lemma (\Cref{lem:schur}) applied
componentwise: part (1) kills cross terms between inequivalent irreps, and
part (2) forces the action on each copy of a given irrep within an isotypic
block to be a shared scalar, leaving only the multiplicity-space map
$A_\pi$---an $m_\pi' \times m_\pi$ matrix, independent of $\dim \pi$---as the
free parameter. This is the entire content of the parameter-count reduction
already flagged at the end of \Cref{subsubsec:schurs-lemma}: a general
linear map $V \to V'$ has $\dim V \cdot \dim V'$ parameters, while the most
general \emph{equivariant} one has only
$\sum_\pi m_\pi m_\pi'$---independent of the (typically much larger)
dimensions $\dim \pi$ of the irreps themselves.

\begin{Example}[Isotypic decomposition constraining architecture]
\label{ex:isotypic-architecture}
\Cref{prop:schur-blocks} is the mechanism, in one form or
another, behind each of the following constructions:
\begin{itemize}
    \item \textbf{Steerable CNNs.} Feature fields on $\mathbb{R}^2$
    equivariant under $E(2) = \mathbb{R}^2 \rtimes O(2)$
    (\Cref{ex:semidirect-products}) are built by assigning each
    feature channel to an isotypic component of a chosen irrep of the
    compact factor $O(2)$ (or a finite subgroup thereof), so that admissible
    convolution kernels are exactly those satisfying the block-diagonal
    constraint of \Cref{prop:schur-blocks}, reducing kernel
    design to solving a small linear system per pair of irreps rather than
    over the full spatial kernel~\citep{weiler2019anips}.
    \item \textbf{Multiplicity-space architectures.} More recent
    equivariant architectures push the isotypic viewpoint further by
    treating the multiplicity spaces $\mathbb{C}^{m_\pi}$ themselves as the
    sole locus of learnable, unconstrained linear structure at each layer,
    with all cross-irrep mixing explicitly disallowed by
    \Cref{prop:schur-blocks}, as in the architecture studied
    in~\citep{xu2022p3icml}.
    \item \textbf{Universality and expressivity limits.} Because
    \Cref{prop:schur-blocks} caps the free parameters of an
    equivariant layer at $\sum_\pi m_\pi m_\pi'$, isotypic decomposition
    also governs what an equivariant network \emph{cannot} express;
    quantifying this gap---how much approximation power is sacrificed by
    imposing the isotypic block constraint relative to an unconstrained
    network of the same width---is the subject of the universal
    approximation analysis in~\citep{yarotsky2018r}.
    \item \textbf{Beyond compact symmetry.} Isotypic decomposition, exactly
    as stated, requires the Peter--Weyl guarantee of
    \Cref{thm:peter-weyl} and hence compactness of $G$
    (\Cref{subsubsec:noncompact-representations}); extending an isotypic-style
    block constraint to symmetry groups with a non-compact factor (e.g.\ the
    translation part of $E(n)$, or scaling symmetries
    $\mathbb{R}_{>0}$) requires replacing the discrete sum over
    $\widehat{G}$ with the direct-integral / Plancherel-measure formalism
    of \Cref{subsubsec:noncompact-representations}, an approach pursued
    in~\citep{gibson2024}.
\end{itemize}
\end{Example}

\subsection{Harmonic Analysis on Groups}
\label{subsec:harmonic-analysis}

The Peter--Weyl decomposition of \Cref{subsubsec:peter-weyl} is, properly
understood, a generalization of the Fourier transform, and the isotypic
constructions of \Cref{sec:isotypic-decomposition} are best understood
as \emph{spectral} constraints in this generalized sense. This subsection
makes that identification explicit and draws out its consequences for
learning.

\subsubsection{Generalities}
\label{subsubsec:harmonic-analysis-generalities}

\begin{Definition}[Fourier transform on a compact group]
\label{def:group-fourier-transform}
Let $G$ be a compact group with normalized Haar measure $\mu$
(\Cref{prop:compact-haar-finite}), and let $\widehat{G}$ be its
unitary dual. For $f \in L^2(G,\mu)$, the \emph{(non-commutative) Fourier
transform} of $f$ is the assignment, to each irrep $\pi \in \widehat{G}$
(realized concretely as $\pi : G \to U(d_\pi)$ on $\mathbb{C}^{d_\pi}$), of
the $d_\pi \times d_\pi$ matrix
\[
    \widehat{f}(\pi) \;:=\; \int_G f(g)\, \pi(g)^{-1} \, d\mu(g) \;\in\; \mathbb{C}^{d_\pi \times d_\pi}.
\]
The Peter--Weyl theorem (\Cref{thm:peter-weyl}) is exactly the
statement that $f \mapsto \big(\widehat{f}(\pi)\big)_{\pi \in \widehat{G}}$
is an isometric isomorphism $L^2(G,\mu) \to \bigoplus_\pi \mathbb{C}^{d_\pi
\times d_\pi}$ (up to normalization by $d_\pi$), with inversion formula
$f(g) = \sum_{\pi \in \widehat{G}} d_\pi \,\operatorname{tr}\bigl(\widehat{f}(\pi)\,\pi(g)\bigr)$.
For $G$ abelian, every irrep is $1$-dimensional (Schur's lemma,
\Cref{lem:schur}(2), applied to $\rho(h)$ commuting with $\rho(g)$ for
all $g,h$ forces every $\rho(g)$ to be scalar), $\widehat{G}$ recovers the
usual group of characters $\chi : G \to U(1)$, and
\Cref{def:group-fourier-transform} reduces exactly to the
classical Fourier transform: on $G = \mathbb{Z}/n\mathbb{Z}$ it is the
discrete Fourier transform, and on $G = \mathbb{R}$ (non-compact, treated
via the direct-integral formalism of
\Cref{subsubsec:noncompact-representations} rather than a discrete sum) it
is the ordinary Fourier transform on the line.
\end{Definition}

The essential feature that makes this transform useful, exactly as in the
classical case, is that it diagonalizes convolution.

\begin{Proposition}[Convolution theorem]
\label{prop:convolution-theorem}
For $f_1, f_2 \in L^2(G,\mu)$, define the \emph{convolution}
\[
    (f_1 * f_2)(g) \;:=\; \int_G f_1(h)\, f_2(h^{-1}g)\, d\mu(h).
\]
Then $\widehat{f_1 * f_2}(\pi) = \widehat{f_1}(\pi)\, \widehat{f_2}(\pi)$ for
every $\pi \in \widehat{G}$: convolution on $G$ becomes ordinary matrix
multiplication after Fourier transform.
\end{Proposition}

Convolution as defined here is exactly the operation with respect to which
$G$ acts on $L^2(G,\mu)$ by left translation, $(g \cdot f)(x) = f(g^{-1}x)$,
a special case of the \emph{regular representation} of
\Cref{ex:representations-recap} carried over from $\mathbb{C}[G]$ to
$L^2(G,\mu)$; \Cref{prop:convolution-theorem} therefore says that
the regular representation is, block by block, exactly the isotypic
decomposition of \Cref{sec:isotypic-decomposition} applied to itself,
with each irrep $\pi$ appearing in $L^2(G,\mu)$ with multiplicity
$m_\pi = d_\pi$ (part (3) of \Cref{thm:peter-weyl}).

For $G = SO(3)$ acting on the sphere $S^2 = SO(3)/SO(2)$, this specializes
to classical spherical harmonic analysis: the irreps of $SO(3)$ are indexed
by $\ell = 0,1,2,\dots$ with $d_\ell = 2\ell+1$, the isotypic components of
$L^2(S^2)$ are exactly the spans of the degree-$\ell$ spherical harmonics,
and \Cref{prop:convolution-theorem} is the statement, familiar
from physics, that convolution with a rotationally symmetric kernel acts
diagonally, by a single scalar per degree $\ell$, on the spherical harmonic
expansion of a function. A systematic treatment of this material at the
level of generality used here is given in Folland's
\emph{A Course in Abstract Harmonic Analysis}~\citep{folland2016}.

\subsubsection{Relation to learning}
\label{subsubsec:harmonic-analysis-learning}

The relevance of this formalism to the learning constructions surveyed in
\Cref{sec:isotypic-decomposition} is that a $G$-equivariant linear
layer is, in these terms, nothing other than a \emph{spectral filter}.

\begin{itemize}
    \item \textbf{Group convolution as the universal equivariant linear
    map.} Any $G$-equivariant linear map between two copies of the regular
    representation $L^2(G,\mu)$ is necessarily convolution by some fixed
    kernel $\kappa \in L^2(G,\mu)$, $f \mapsto \kappa * f$; this is a special
    case of \Cref{prop:schur-blocks} in which the multiplicity
    spaces $\mathbb{C}^{m_\pi}$ have dimension $d_\pi$ and $A_\pi$ is
    determined by $\widehat{\kappa}(\pi)$. This is precisely the
    group-convolutional layer of~\citep{cohen2014}: replacing ordinary
    translation-convolution by convolution over a general (typically finite
    or compact) symmetry group $G$ is, by
    \Cref{prop:convolution-theorem}, automatically equivariant,
    and the isotypic block-diagonal parameter count of
    \Cref{sec:isotypic-decomposition} is recovered by passing to the
    Fourier side, where the layer is literally a block-diagonal matrix
    $\bigl(\widehat{\kappa}(\pi)\bigr)_{\pi}$.
    \item \textbf{Learning in the spectral domain.} Because
    \Cref{prop:convolution-theorem} makes the layer diagonal on the Fourier
    side, architectures for compact $G$ frequently
    parametrize their learnable kernels directly there, storing one small
    matrix $\widehat{\kappa}(\pi)$ per
    isotypic block---at the cost of a forward/inverse group Fourier
    transform (or, on $S^2$, a spherical harmonic transform) at each layer
    boundary. Spherical CNNs \citep{cohen2018} implement exactly this,
    using a generalized FFT on $SO(3)$, and the Clebsch--Gordan networks
    of \citet{kondor2018p3icml} go further by never leaving the spectral
    domain at all, performing the nonlinearity itself on Fourier
    coefficients so that no transform is needed between layers.
    \item \textbf{Non-compact groups again lose this structure.} As
    flagged in \Cref{subsubsec:noncompact-representations} and revisited in
    \Cref{ex:isotypic-architecture}, once $G$ has a non-compact
    factor (translations, positive rescalings), $\widehat{G}$ is no longer
    a discrete index set, the sums above become direct integrals against a
    Plancherel measure, and ``one scalar per irrep'' becomes ``one scalar per
    point of a continuous spectral parameter'' (ordinary frequency, in the
    translation case)---recovering exactly classical Fourier analysis on
    $\mathbb{R}^n$ as the $G = \mathbb{R}^n$ special case of this same
    formalism, but without the finite, enumerable spectral structure that
    made the compact case so directly exploitable for parameter reduction.
\end{itemize}

Harmonic analysis on groups is thus the unifying language behind
\Cref{sec:isotypic-decomposition}: an isotypic decomposition
\emph{is} a Fourier decomposition with respect to $G$, an equivariant linear
layer \emph{is} a (possibly group-)convolution, and Schur's lemma
(\Cref{lem:schur}) is, in this dictionary, precisely the statement that
convolution operators are diagonalized by the group Fourier transform.


% Homogeneous Spaces cohen2019anips

\subsection{Homogeneous Spaces}
\label{subsec:homogeneous-spaces}

The sphere $S^2 = SO(3)/SO(2)$ used to illustrate harmonic analysis in
\Cref{subsubsec:harmonic-analysis-generalities} is an instance of a more
general and recurring pattern: many of the spaces on which learning problems
are naturally posed are not themselves groups, but arise as a group modulo
a subgroup of ``directions one is willing to forget.'' This subsection makes
that pattern precise and connects it back to the orbit--stabilizer picture
of \Cref{subsec:group-actions}.

\subsubsection{Definition and basic examples}
\label{subsubsec:homogeneous-spaces-def}

\begin{Definition}[Homogeneous space]
\label{def:homogeneous-space}
A \emph{homogeneous space} for a group $G$ is a set (or manifold) $\mathcal{X}$
equipped with a \emph{transitive} action of $G$
(\Cref{def:free-transitive}). Given a point $x_0 \in \mathcal{X}$ with
stabilizer $H := G_{x_0} \leq G$ (\Cref{def:orbit-stabilizer}),
transitivity means $\mathcal{X}$ is the single orbit $G \cdot x_0$, so the
Orbit--Stabilizer Theorem (\Cref{thm:orbit-stabilizer}) upgrades from
a statement about one orbit inside a possibly larger $\mathcal{X}$ to a statement
identifying \emph{all} of $\mathcal{X}$:
\[
    \mathcal{X} \;\cong\; G / H
\]
as $G$-sets (and, when $G$ is a Lie group acting smoothly and properly, as
smooth manifolds, per the remark closing
\Cref{subsubsec:lie-groups}), via $gH \mapsto g \cdot x_0$. Different choices
of basepoint $x_0$ give conjugate stabilizers and hence isomorphic, but not
literally identical, coset descriptions $G/H$.
\end{Definition}

Homogeneous spaces thus sit at exactly one further level of generality above
the $G$-torsors of \Cref{def:free-transitive}: a torsor is the
special case $H = \{e\}$, i.e.\ a homogeneous space with trivial stabilizer,
so that $\mathcal{X} \cong G$ itself once a basepoint is fixed. Allowing $H$ nontrivial
is what lets $\mathcal{X}$ be genuinely smaller than $G$.

\begin{Example}[Homogeneous spaces already in use]
\label{ex:homogeneous-spaces}
\begin{itemize}
    \item $S^{n-1} = SO(n)/SO(n-1)$: the sphere, with $SO(n)$ acting
    transitively by rotation and $SO(n-1)$ the stabilizer of a fixed pole
    (rotations fixing a chosen axis)---already used implicitly for $n=3$ in
    \Cref{subsubsec:harmonic-analysis-generalities}.
    \item $\mathbb{R}^n = E(n) / O(n)$: Euclidean space itself, with the
    full Euclidean group $E(n)$ of \Cref{subsubsec:semidirect-products}
    acting transitively by rigid motion and $O(n)$ the stabilizer of the
    origin---exhibiting $\mathbb{R}^n$ as a homogeneous, but not free
    (hence not torsor), space for $E(n)$, in contrast to the free
    translation action of $\mathbb{R}^n \leq E(n)$ alone.
    \item $G/G = \{pt\}$ and $G/\{e\} = G$ are the two trivial extremes,
    recovering the single-orbit and torsor cases respectively as special
    instances of \Cref{def:homogeneous-space}.
    \item More generally, \emph{Grassmannians} $\mathrm{Gr}(k,n) =
    O(n)/(O(k)\times O(n-k))$, parametrizing $k$-dimensional subspaces of
    $\mathbb{R}^n$, arise the same way, with the stabilizer of a fixed
    $k$-plane splitting as a product acting separately on the plane and its
    orthogonal complement.
\end{itemize}
\end{Example}

\subsubsection{Role in learning}
\label{subsubsec:homogeneous-spaces-learning}

Homogeneous spaces matter for learning precisely where the isotypic and
harmonic-analytic machinery of
\Cref{sec:isotypic-decomposition}--\Cref{subsec:harmonic-analysis}
needs a \emph{domain} to operate on that is not the group $G$ itself, but a
smaller space $\mathcal{X} = G/H$ on which $G$ still acts transitively---this is
exactly the situation for data living on a sphere, a Grassmannian, or any
other space with less symmetry-breaking structure than the full group.

\begin{itemize}
    \item \textbf{Functions on $G/H$ as an $H$-invariant subspace of
    functions on $G$.} Pulling back along the quotient map $G \to G/H$
    identifies $L^2(G/H,\mu)$ with the subspace of $L^2(G,\mu)$ consisting
    of functions invariant under right translation by $H$; under the
    Peter--Weyl decomposition of
    \Cref{def:group-fourier-transform}, this subspace is itself
    isotypic-block-structured, being spanned by exactly those matrix
    coefficients $\pi(g)_{ij}$ for which the corresponding basis vector is
    fixed by $H$---so the harmonic analysis of $G/H$ reduces to bookkeeping
    which irreps of $G$ contain a nonzero $H$-fixed vector, and with what
    multiplicity.
    \item \textbf{Equivariant networks on homogeneous spaces.} This is
    precisely the perspective developed systematically
    in~\citet{cohen2019anips}, which constructs the general theory of
    convolutional layers equivariant under a group $G$ acting on data
    supported on a homogeneous space $\mathcal{X} = G/H$ (subsuming, as special cases,
    planar CNNs with $G = \mathbb{Z}^2 \rtimes$(point group), $H$ trivial,
    and spherical CNNs with $G = SO(3)$, $H = SO(2)$ as in
    \Cref{ex:homogeneous-spaces}), showing that the admissible
    linear, equivariant maps between feature fields on $G/H$ are again
    governed by an isotypic/Schur's-lemma constraint
    (\Cref{prop:schur-blocks}) once the group Fourier transform
    of \Cref{subsubsec:harmonic-analysis-generalities} is restricted to the
    $H$-fixed subspace described above.
    \item \textbf{Non-compact base spaces.} When $G$ has a non-compact
    factor---e.g.\ $\mathcal{X} = \mathbb{R}^n = E(n)/O(n)$ itself---the same
    quotient construction applies formally, but the underlying harmonic
    analysis reverts to the direct-integral, Plancherel-measure regime of
    \Cref{subsubsec:noncompact-representations} rather than the discrete
    isotypic sum, exactly mirroring the compact/non-compact split already
    emphasized throughout this section.
\end{itemize}

The upshot is that ``equivariant network for symmetry group $G$'' and
``network operating on the homogeneous space $\mathcal{X} = G/H$'' are, in this
framework, two descriptions of the same underlying construction: the choice
of $\mathcal{X}$ fixes the stabilizer $H$, the stabilizer fixes which isotypic
components of $L^2(G)$ survive restriction to $\mathcal{X}$, and Schur's lemma then
fixes, exactly as in \Cref{sec:isotypic-decomposition}, the space of
admissible equivariant linear layers on $\mathcal{X}$.


% Advanced notions
% Structures beyond groups (no inverse: semigroup)
% Tensor algebra
% spherical harmonics (used a lot in Euclidean neural networks)
% Lie algebra, Exponential and logarithmic map, Infinitesimal generators
% Manifold/Gauge Theory: fiber bundles, Principal bundle, section, steerable kernels, local symmetries, free action
% Symmetry breaking

\subsection{Advanced Notions}
\label{sec:advanced-notions}

We close this survey of the algebraic backbone with six topics that sit one
level beyond the material developed so far: structures that weaken the group
axioms, the analytic and infinitesimal apparatus attached to a Lie group,
and the geometric (bundle-theoretic) language required once symmetry is
allowed to vary from point to point rather than acting globally. Each
connects directly to a construction already introduced.

\subsubsection{Structures Beyond Groups: Semigroups and Monoids}
\label{subsec:beyond-groups}

Relaxing the group axioms of \Cref{def:group} produces weaker, but
still useful, algebraic structures---most importantly by dropping (G4).

\begin{Definition}[Semigroup and monoid]
\label{def:semigroup-monoid}
A \emph{semigroup} is a pair $(S,\cdot)$ satisfying only closure (G1) and
associativity (G2): a set with an associative binary operation, but no
requirement of an identity or of inverses. A \emph{monoid} is a semigroup
additionally satisfying the identity axiom (G3). A group is thus a monoid in
which every element happens to be invertible.
\end{Definition}

\begin{Example}[Semigroups and monoids arising from learning constructions]
\label{ex:semigroups-in-learning}
\begin{itemize}
    \item The natural numbers $(\mathbb{N}, +)$ form a monoid but not a
    group: $0$ is an identity, but no positive integer has an additive
    inverse in $\mathbb{N}$.
    \item Square matrices $\operatorname{Mat}_n(\mathbb{R})$ under multiplication
    form a monoid (identity $I$), of which the invertible matrices
    $GL(n,\mathbb{R})$ of \Cref{ex:groups} are precisely the
    invertible elements---i.e.\ $GL(n,\mathbb{R})$ is the group of units of
    the monoid $\operatorname{Mat}_n(\mathbb{R})$. Layers of a network built from
    non-invertible linear maps (e.g.\ any layer that reduces dimension) thus
    naturally compose as a monoid, not a group, under composition.
    \item The rescaling symmetry $\mathbb{R}_{>0}$ of
    \Cref{subsubsec:lie-groups} is a group, but its closure
    $\mathbb{R}_{\geq 0}$ under multiplication (allowing the scale factor to
    reach zero) is only a monoid: $0$ has no multiplicative inverse. This
    matters directly for ReLU-type rescaling symmetries, where a weight
    vector is allowed to degenerate to exactly zero along a training
    trajectory (a ``dead unit''); the exact symmetry group
    $\mathbb{R}_{>0}$ acting on nonzero weights extends only to a
    monoid action once the zero locus---a lower-dimensional stratum with
    strictly smaller stabilizer-orbit structure than
    \Cref{thm:orbit-stabilizer} describes for a genuine group
    action---is included.
    \item The set $\operatorname{End}(V)$ of all linear self-maps of a vector
    space $V$ is a monoid under composition, and every representation
    $\rho: G \to GL(V) \subset \operatorname{End}(V)$
    (\Cref{def:linear-representation}) is, forgetting
    invertibility, in particular a monoid homomorphism $G \to
    \operatorname{End}(V)$; this is the observation that lets representation
    theory be extended, with care, to monoid and semigroup actions where
    needed.
\end{itemize}
\end{Example}

The loss of inverses is consequential: orbits of a semigroup action need not
partition the underlying set into disjoint equivalence classes (since
$x \sim y$ via $s \cdot x = y$ no longer implies $y \sim x$), so neither
\Cref{thm:orbit-stabilizer} nor the isotypic machinery of
\Cref{sec:isotypic-decomposition} extends verbatim; monoid and
semigroup actions are consequently treated as a genuinely separate,
order-theoretic (rather than group-theoretic) subject in general.

\subsubsection{Tensor Algebra Revisited: Tensor Fields}
\label{subsec:tensor-fields}

The tensor algebra of a single representation
(\Cref{def:tensor-operations}) globalizes, once $G$ acts on a
manifold $\mathcal{X}$ rather than a single vector space, to a \emph{tensor field}: a
choice of tensor at every point, transforming consistently under the action.

\begin{Definition}[Tensor field]
\label{def:tensor-field}
Let $\mathcal{X}$ be a manifold with a smooth action of a Lie group $G$
(\Cref{def:group-action}), and fix a representation
$\rho : G \to GL(V)$. A \emph{$V$-valued tensor field of type $\rho$} is a
smooth assignment $x \mapsto \tau(x) \in V$ (or, more generally, a section of
the vector bundle $\mathcal{X} \times_\rho V$ associated to $\rho$, in the sense made
precise in \Cref{subsec:manifold-gauge-theory} below) such that
\[
    \tau(g \cdot x) \;=\; \rho(g)\, \tau(x) \qquad \text{for all } g \in G,\ x \in \mathcal{X}.
\]
The tangent, cotangent, and higher tensor bundles $T\mathcal{X}$, $T^*\mathcal{X}$,
$T^{\otimes k}\mathcal{X}$ are the special cases in which $\rho$ is the (co)tangent
representation induced by $G$ acting on $\mathcal{X}$ by diffeomorphism, so that
ordinary tensor fields of differential geometry are recovered as the case
$V = T_x^{\otimes k}\mathcal{X}$.
\end{Definition}

This is the object directly manipulated in equivariant architectures that
output more than a scalar per point: a steerable CNN feature map
(\Cref{ex:isotypic-architecture}) is exactly a tensor field of type
$\rho$ for $\rho$ an isotypic sum of irreps of the stabilizer, and the
transformation law above is the pointwise instance of the global
equivariance constraint enforced by \Cref{prop:schur-blocks}.

\subsubsection{Spherical Harmonics}
\label{subsec:spherical-harmonics}

Spherical harmonics were introduced implicitly in
\Cref{subsubsec:harmonic-analysis-generalities} as the isotypic components
of $L^2(S^2)$; we record their defining properties explicitly, since they
are used pervasively, by name, throughout the Euclidean-equivariant network
literature.

\begin{Definition}[Spherical harmonics]
\label{def:spherical-harmonics}
For $\ell = 0, 1, 2, \dots$, the space of \emph{degree-$\ell$ spherical
harmonics} $\mathcal{H}_\ell \subset L^2(S^2)$ is the space of restrictions
to $S^2 \subset \mathbb{R}^3$ of harmonic, degree-$\ell$ homogeneous
polynomials on $\mathbb{R}^3$ (i.e.\ $p$ with $\Delta p = 0$ and
$p(\lambda x) = \lambda^\ell p(x)$). It has dimension $2\ell+1$, is invariant
under the $SO(3)$-action of \Cref{ex:homogeneous-spaces}, and is
precisely the irrep of $SO(3)$ of ``degree $\ell$'' appearing in the
Peter--Weyl decomposition
\[
    L^2(S^2) \;=\; \bigoplus_{\ell = 0}^{\infty} \mathcal{H}_\ell,
\]
a special case of \Cref{def:group-fourier-transform} for
$G = SO(3)$, $X = S^2 = SO(3)/SO(2)$.
\end{Definition}

Spherical harmonics enter Euclidean-equivariant network design because they
are simultaneously (a) an explicit, orthonormal basis diagonalizing
rotation-equivariant linear maps on functions of direction, per
\Cref{prop:convolution-theorem}, and (b) the natural target
representation type for features attached to \emph{directions} (bond
vectors, edge orientations, point-cloud relative positions) in 3D geometric
learning problems---so that a network layer producing, say, a degree-$\ell$
feature is by construction guaranteed to transform correctly under any
ambient $SO(3)$ or $E(3)$ rotation, without the transformation law needing
to be learned or enforced by a penalty term. Tensor-field networks and the
spherical-CNN/Clebsch--Gordan constructions of \citet{kondor2018p3icml} use
the Clebsch--Gordan coefficients of \Cref{subsubsec:tensor-algebra} exactly
to combine features of different degrees $\ell_1, \ell_2$ into new features
of degree $\ell_3$ while staying inside this basis throughout.

\subsubsection{Lie Algebras, Exponential and Logarithmic Maps, Infinitesimal Generators}
\label{subsec:lie-algebras}

We now make precise the linearization $\mathfrak{g} = T_eG$ of a Lie group
$G$ mentioned in passing in \Cref{def:lie-group}.

\begin{Definition}[Lie algebra]
\label{def:lie-algebra}
The \emph{Lie algebra} of a Lie group $G$ is the tangent space
$\mathfrak{g} := T_e G$ at the identity, equipped with the \emph{Lie
bracket} $[\cdot,\cdot] : \mathfrak{g} \times \mathfrak{g} \to \mathfrak{g}$
induced by the failure of left- and right-translated vector fields to
commute (equivalently, the derivative at $e$ of the commutator map
$(g,h) \mapsto ghg^{-1}h^{-1}$ in both variables). The bracket is bilinear,
antisymmetric ($[X,Y] = -[Y,X]$), and satisfies the Jacobi identity
$[X,[Y,Z]] + [Y,[Z,X]] + [Z,[X,Y]] = 0$.
\end{Definition}

\begin{Example}[Matrix Lie algebras]
\label{ex:matrix-lie-algebras}
For $G \leq GL(n,\mathbb{R})$ a matrix Lie group, $\mathfrak{g}$ is
concretely realized as a linear subspace of matrices, with
$[X,Y] = XY - YX$ the ordinary matrix commutator:
$\mathfrak{gl}(n,\mathbb{R}) = \operatorname{Mat}_n(\mathbb{R})$ for
$G = GL(n,\mathbb{R})$; $\mathfrak{so}(n) = \{X : X^\top = -X\}$
(antisymmetric matrices) for $G = SO(n)$, obtained by differentiating
$A(t)^\top A(t) = I$ at $t=0$ along a curve $A(t)$ with $A(0)=I$; and
$\mathfrak{so}(1,n) = \{X : X^\top \eta = -\eta X\}$ for the Lorentz group
$O(1,n)$ of \Cref{subsubsec:lie-groups}, with $\eta = \operatorname{diag}(-1,1,\dots,1)$
the Minkowski metric, by the identical computation.
\end{Example}

\begin{Definition}[Exponential and logarithmic map]
\label{def:exp-log-map}
The \emph{exponential map} $\exp : \mathfrak{g} \to G$ sends $X \in
\mathfrak{g}$ to $\gamma_X(1)$, where $\gamma_X : \mathbb{R} \to G$ is the
unique one-parameter subgroup ($\gamma_X(s+t) = \gamma_X(s)\gamma_X(t)$)
with $\gamma_X'(0) = X$; for matrix groups, $\exp(X) = \sum_{k \geq 0}
X^k/k!$, the ordinary matrix exponential. $\exp$ is a local
diffeomorphism near $0 \in \mathfrak{g}$ (and near $e \in G$), with local
inverse the \emph{logarithmic map} $\log : G \to \mathfrak{g}$ near $e$; for
$G$ connected and compact, $\exp$ is in fact globally surjective (though
generally not injective), while for non-compact $G$---consistent with the
recurring compact/non-compact split of
\Cref{subsubsec:noncompact-representations}---surjectivity of $\exp$ can
fail (e.g.\ not every element of $SL(2,\mathbb{R})$ is $\exp$ of a single
element of $\mathfrak{sl}(2,\mathbb{R})$).
\end{Definition}

\begin{Definition}[Infinitesimal generators]
\label{def:infinitesimal-generators}
Given a representation $\rho : G \to GL(V)$
(\Cref{def:linear-representation}) of a Lie group $G$, its
derivative at the identity, $d\rho_e : \mathfrak{g} \to \mathfrak{gl}(V) =
\operatorname{End}(V)$, is a Lie algebra representation (a linear map preserving
the bracket, $d\rho_e[X,Y] = [d\rho_e X, d\rho_e Y]$), and the elements
$d\rho_e(X)$ for $X \in \mathfrak{g}$ are the \emph{infinitesimal
generators} of $\rho$. They recover the group representation by
exponentiation, $\rho(\exp X) = \exp(d\rho_e(X))$, so that on a connected
Lie group, $\rho$ is entirely determined by the (typically far more
tractable, since linear rather than group-theoretic) data of $d\rho_e$.
\end{Definition}

This linearization is what makes symmetry constraints tractable
computationally: enforcing exact equivariance under the full group $G$ is,
by \Cref{def:infinitesimal-generators}, equivalent (at least
locally, and globally when $G$ is connected) to enforcing equivariance under
the finitely many infinitesimal generators of $\mathfrak{g}$---a Lie
algebra, hence a linear, condition. This is the mechanism behind
Lie-algebra-based approaches to learning approximate or exact symmetries
directly from data, where the generators $d\rho_e(X)$ themselves (rather
than the full nonlinear group action) are the learned or regularized
objects.

\subsubsection{Manifolds and Gauge Theory: Fiber Bundles, Principal Bundles, Local Symmetries}
\label{subsec:manifold-gauge-theory}

The final piece of geometric machinery is needed once symmetry is allowed to
vary smoothly from point to point of a base space, rather than acting
globally and uniformly as in every construction so far.

\begin{Definition}[Fiber bundle]
\label{def:fiber-bundle}
A \emph{fiber bundle} is a surjective smooth map $\pi : E \to \mathcal{X}$ between
manifolds such that every $x \in \mathcal{X}$ has a neighborhood $U$ with
$\pi^{-1}(U) \cong U \times F$ (a \emph{local trivialization}) for a fixed
model fiber $F$, compatibly with $\pi$. A \emph{section} of the bundle is a
smooth map $s : \mathcal{X} \to E$ with $\pi \circ s = \operatorname{id}_X$, i.e.\ a smooth
choice of one point of the fiber $\pi^{-1}(x)$ for each $x$; the tensor
fields of \Cref{def:tensor-field} are exactly sections of the
associated bundle $\mathcal{X} \times_\rho V$ mentioned there.
\end{Definition}

\begin{Definition}[Principal bundle]
\label{def:principal-bundle}
A \emph{principal $G$-bundle} is a fiber bundle $\pi : P \to \mathcal{X}$ with model
fiber $G$ itself, equipped with a smooth, fiberwise, free and transitive
right action of $G$ on $P$ (\Cref{def:free-transitive} applied
fiberwise: each fiber $\pi^{-1}(x)$ is a $G$-torsor). Local
trivializations correspond to local choices of section (a local ``reference
frame,'' or \emph{gauge}), and two local sections over overlapping
neighborhoods differ by a smooth \emph{transition function} $g_{\alpha\beta}
: U_\alpha \cap U_\beta \to G$---the mathematical formalization of a
\emph{local (gauge) symmetry}: the group $G$ acting freely and transitively
on each fiber, as opposed to acting globally on all of $\mathcal{X}$ at once as in
every construction of \Cref{sec:group-theory-backbone}--\Cref{subsec:homogeneous-spaces}.
\end{Definition}

\begin{Example}[Principal bundles in geometric learning]
\label{ex:principal-bundles-learning}
\begin{itemize}
    \item The \emph{frame bundle} of a manifold $\mathcal{X}$, whose fiber over $x$ is
    the set of ordered bases of $T_x\mathcal{X}$, is a principal $GL(n,\mathbb{R})$-bundle
    (or principal $O(n)$-bundle, once restricted to orthonormal frames); a
    choice of local orthonormal frame field is exactly a local gauge, and
    changing frame at each point by a rotation is the geometric meaning of
    the local $O(n)$ symmetry in this setting.
    \item \textbf{Gauge-equivariant / steerable convolution on a curved
    manifold.} Where \Cref{sec:isotypic-decomposition} constrained
    linear maps equivariant under a \emph{global} symmetry $G$ acting on
    flat $\mathbb{R}^n$ or a homogeneous space $G/H$, gauge-equivariant CNNs
    replace the global $G$-action by the structure group of a principal
    bundle over a general (possibly curved, possibly non-homogeneous)
    manifold $\mathcal{X}$, and require that a convolution kernel, transported between
    two local gauges related by a transition function $g_{\alpha\beta}$, be
    equivariant under the fiberwise action of $G$ in exactly the sense of
    \Cref{prop:schur-blocks}, applied now pointwise rather than
    globally---this is the \emph{steerable kernel constraint} that
    generalizes the flat steerable-CNN condition of~\citep{weiler2019anips}
    from a homogeneous space (\Cref{subsec:homogeneous-spaces}) to an
    arbitrary manifold equipped only with local, rather than global,
    symmetry.
    \item Associated vector bundles $\mathcal{X} \times_\rho V$
    (\Cref{def:fiber-bundle}) built from a representation $\rho$
    of the structure group $G$ of a principal bundle are exactly the
    bundles whose sections are the tensor fields of
    \Cref{def:tensor-field}; a gauge transformation acts on such a
    section by $\rho$ applied fiberwise, and the free, transitive fiberwise
    action of \Cref{def:principal-bundle} guarantees this action
    is well defined independent of the (arbitrary) choice of local frame.
\end{itemize}
\end{Example}

The key conceptual shift, relative to every earlier construction in this
survey, is captured by the word \emph{local}: \Cref{subsec:group-actions}
through \Cref{subsec:homogeneous-spaces} concerned a single group acting
once, globally, on an entire space $\mathcal{X}$; gauge theory concerns a group acting
independently, and freely, on each fiber over each point of $\mathcal{X}$
separately, with the transition functions $g_{\alpha\beta}$ recording how
these independent local actions are glued together consistently across
overlapping patches. Equivariance in this setting is consequently a
\emph{local}, gauge-covariant condition (transforming correctly under
independent reparametrization of each patch) rather than a single global
group-homomorphism condition as in \Cref{def:linear-representation}.

\subsubsection{Symmetry Breaking}
\label{subsec:symmetry-breaking}

We close with the phenomenon that determines when the symmetry-respecting
machinery of this entire section applies to a given problem \emph{at all}.

\begin{Definition}[Symmetry breaking]
\label{def:symmetry-breaking}
Let $G$ act on a set (or space) $\mathcal{Y}$ of admissible states, models, or ground
truths for a given problem. \emph{Symmetry breaking} occurs when the actual
object of interest $y_0 \in \mathcal{Y}$ (a physical ground state, a trained network,
a data distribution) has stabilizer $G_{y_0} \subsetneq G$
(\Cref{def:orbit-stabilizer}) strictly smaller than the full
symmetry group of the ambient problem---equivalently, the orbit
$G \cdot y_0$ is nontrivial, so distinct elements of $G$ produce genuinely
distinct, but equally valid, instances $g \cdot y_0$.
\end{Definition}

\begin{Example}[Symmetry breaking already implicit in this survey]\label{ex:symmetry-breaking}
\phantom{u}
\begin{itemize}
    \item \textbf{Weight-space redundancy versus a trained network.} The
    full permutation/rescaling group $S_n \ltimes \mathbb{R}_{>0}^n$
    (\Cref{ex:semidirect-products}) acts on the space of all
    parameterizations of a hidden layer with trivial stabilizer generically
    (a free action, per \Cref{def:free-transitive}, away from
    degenerate configurations); a specific trained weight vector $w_0$
    \emph{breaks} this symmetry down to the trivial subgroup, and the orbit
    $G \cdot w_0$ is exactly the set of functionally identical
    reparameterizations quotiented out in
    \Cref{sec:parameter-space-symmetries}. Here symmetry breaking is
    generic and unwanted---the entire motivation for passing to the
    quotient $\mathcal{Y}/G$ in the first place.
    \item \textbf{Physical symmetry breaking.} In the physical setting that
    originally motivates the term, a Hamiltonian invariant under the full
    Euclidean or Lorentz group $E(n)$ or $O(1,n)$
    (\Cref{subsubsec:lie-groups}) can admit a ground state $y_0$ invariant
    only under a strict subgroup $H \leq G$ (e.g.\ a crystal breaking
    continuous translation symmetry down to a discrete lattice subgroup,
    connecting back to \Cref{ex:discrete-groups}); the set of
    physically equivalent ground states is then the homogeneous space
    $G/H$ of \Cref{subsec:homogeneous-spaces} itself, with $H$ playing the
    role of the unbroken stabilizer.
    \item \textbf{Symmetry breaking as an architectural design choice.}
    Equivariant architectures built from the isotypic
    (\Cref{sec:isotypic-decomposition}) or gauge-theoretic
    (\Cref{subsec:manifold-gauge-theory}) machinery above are exactly
    calibrated \emph{not} to break a chosen symmetry $G$ at the level of the
    function class; whether the learned network's specific output on a
    specific input further breaks symmetry (e.g.\ classifying a
    rotationally ambiguous image into one class rather than a rotated
    alternative) is then a separate, data-dependent question about the
    orbit of that particular output under the residual symmetry of the
    input, in the sense of \Cref{def:symmetry-breaking} applied
    to $\mathcal{Y} = $ the space of possible outputs.
\end{itemize}
\end{Example}

Symmetry breaking is therefore not a single phenomenon but a lens: applied
to \emph{parameters}, it is the redundancy problem motivating the weight-space
quotients of this survey; applied to \emph{ground states} or \emph{learned
solutions}, it is the mechanism by which a fully symmetric problem
nonetheless admits asymmetric, but equally valid, solutions; and in both
cases it is measured by exactly the same orbit--stabilizer pair
(\Cref{def:orbit-stabilizer}) introduced at the very start of this section, closing the loop back to \Cref{thm:orbit-stabilizer} as the single invariant underlying every construction surveyed here.



\bibliographystyle{plainnat}
\bibliography{../bibliography}

\end{document}
